\section{保羅書信專題}

\subsection{從林前12-14, 羅12, 弗4中看保羅對聖靈恩賜（事奉的恩賜）的教導}


\subsubsection{前言}
舊約中「應許的聖靈」在五旬節的時候應驗。這個預言最先記載在先知約珥書中：「以後， 我要將我的靈澆灌凡有血氣的…我要將我的靈澆灌我的僕人和使女。」（珥2:28,29）。耶穌在最後晚餐時，曾安慰門徒說：「我要求父，父就另外賜給你們一位保惠師 ，叫他永遠與你們同在，就是真理的聖靈」（約14:16,17a）。耶穌復活后的第一天晚上同眾門徒在一起的時候，再次提醒門徒等候父所要差遣的聖靈（約20:21,22，路24:49）。耶穌升天後，五旬節門徒聚集在耶路撒冷禱告的時候被聖靈充滿，「應許的聖靈」得到應驗。 「應許的聖靈」是神賜給每一位相信耶穌名的人的印記（徒2:38）。從五旬節信徒在降下聖靈后就開始說方言的記載，可以知道「聖靈的恩賜」與聖靈同時賜給教會。聖靈隨己意將「聖靈的恩賜」分給各人（林前12:4-11）。這些「聖靈的恩賜」是聖靈所賜予，為了使個人及教會得著益處。下面從保羅的書信中來討論聖靈恩賜的相關教導。
\subsubsection{恩賜的目的及意義}
A. 目的
1. 叫人得益處（林前12:7）
a. 使信徒之間有區別 這裡所說的區別表現在不同方面。 1) 言語上：智慧的言語，知識的言語，說方言，翻譯方言和先知的預言。 2) 甚至人可以具有超越人類能力的作為，比如可以醫病，行異能和辨別諸靈等。 3) 工作的超人能力，比如超自然的信心，服事，治理。 b. 使信徒之間彼此相顧 不同的恩賜使得信徒成爲基督身上不同的肢體彼此配搭服事（林前12:12-27）。 1) 保羅用身體的肢體做比喻，解釋我們每個人所得的恩賜不一樣，所以我們可以成為基督身上不同的肢體（林前12:12），因為我們不同，所以我們可以彼此相顧（林前12:25）。 2) 不同的恩賜使得信徒彼此不可缺少，尤其是軟弱的，更是不可缺少的（林前12:22）。 3) 不同的恩賜使得信徒彼此補足，神要用榮耀的肢體去補足不俊美的肢體（林前12:23，24）。 2. 叫人蒙恩典 我們各人蒙恩，都是照基督所量給各人的恩賜使我們蒙恩典（弗4:7）。恩賜是賜給了每一個信徒的（each one of us1）,不是只給了過去的使徒或者今天的牧師，或其他特別的人群。 3. 造就教會（林前14:12） 做先知講道的就是造就的恩賜（林前14:4b）。保羅特別鼓勵信徒要切慕造就教會的恩賜。 4. 造就信徒（林前14:4） 有的恩賜是造就個人的，比如說方言禱告的恩賜（林前14:4a）。
1. Life Application Study Bible（NIV）, (Wheaton, IL: Tyndale House Publishes, Inc., 1991), 2136.
2
B. 意義
1. 恩賜是神的恩典（弗4:7-8）
恩賜不是我們天然的才幹，既然恩賜是恩典，就不是我們靠著自己的努力可以得到的。
2. 使信徒可以做不同的神的工作（羅12:5）
比如說預言的恩賜，治理神的教會的恩賜。
3. 使不同的信徒成為一個身體（羅12:5）
雖然我們恩賜不同，但卻連在一個身體上，一同建造這同一個身體，就是基督的教會。
4. 使信徒彼此不同
恩賜原有分別（林前12:4），不同的恩賜使得信徒成爲基督身上不同的肢體彼此配搭服事（林前12:12-27）。
\subsubsection{領受恩賜的方式 聖靈隨己意運行分配給各人（林前12:11）}
A. 保羅做福音的執事是照神的恩賜，這恩賜是照神運行的大能賜給保羅的（弗3:7）。
B. 領受恩賜的方式
1. 恩賜的來源
恩賜是不同的，但源自同一位聖靈（林前12:4）。
2. 領受恩賜的方式
當人信主後，聖靈要顯在信徒的各人身上（林前12:7），就把對信徒有益處的不同恩賜賜給他們。這種賜給是聖靈主動作的，不是信徒主動要求的，信徒是接受者。
3. 領受什麼恩賜
恩賜有多種，每個信徒身上的恩賜不一樣，這也不是信徒可以要求或決定的，給每個信徒身上什麼樣的恩賜都是聖靈的工作，聖靈是隨己意分給各人的。也就是說每個人身上的恩賜，自己不能決定要和要什麼，或者不要，全是聖靈分配的，人不能拒絕，也不能排斥，人能作的就是把它發揮使用出來，建造教會。
\subsubsection{恩賜的分類}
聖靈是看不見的，聖靈内住在人生命裏也是肉眼看不見的。恩賜是内住的聖靈用肉眼看的見的方式的外在彰顯。恩賜有許多種，按照羅馬書12章和林前12章的分類，共有十五種。每次恩賜彰顯的時候不一定只用到一種，比如翻方言和說方言的恩賜。以下的介紹是基於林前12章和羅馬書12章的分類方法加以討論。
A. 恩賜的種類及分類簡介
恩賜按照彰顯的方式不同，分爲三類：話語的恩賜，異能的恩賜和工作的恩赐。
類別
恩賜的名稱
話語的恩賜
智慧的言語（林前12:8a）
知識的言語（林前12:8b）
先知預言（羅12:6）（林前12:10b）（林前12:28b）
說方言（different kinds of tongues2）（林前12:10d）（林前12:28h）
翻方言（林前12:10e）
教導（羅12:7b）
勸化（羅12:8a）
2. Life Application Study Bible（NIV）, 2081.
3
異能的恩賜
醫病（gifts of healing3）（林前12:9b）（林前12:28e）
行異能（林前12:10a）（林前12:28d）
辨別諸靈（林前12:10c）
工作的恩賜
信心（林前12:9a）
執事Serve4（服事）（羅12:7a）
施捨（羅12:8b）
治理Lead5（羅12:8c）
憐憫（羅12:8d）
B. 話語的恩賜
話語的恩賜包括智慧的言語，知識的言語，先知的預言和講道，使用人間的言語或者是天使的言語說方言（different kinds of tongues6），翻譯方言，教導和勸化。因爲這些恩賜的彰顯都是通過嘴巴，所以都歸類為話語的恩賜。下面用表格來簡述這几種恩賜的説明和使用原則。
恩賜
簡單説明
原則
智慧的言語
不一定每個人都有這個恩賜（林前12:8a）
知識的言語
不一定每個人都有這個恩賜（林前12:8b）
知識的言語需要有愛來搭配（林前13:2）；知識必歸於無有（林前13:8）
先知
在教會中僅次於使徒（林前12:28），可以造就教會，勸勉安慰人（林前14:3）。先知分成以下兩種。
做先知的恩賜是對人說；這個恩賜最終會歸於無有（林前13:8）先知所知道的都是有限的，直到見到主時我們所知道的才能完全（林前13:12）。
1 在教會中說預言的先知 說預言時當照著信心的程度說預言
2 在聚會中講道的先知 做先知講道的人可以明白各樣的奧祕比如保羅（林前13:2）；在聚會中，先知的講道可以造就教會。
聚會時先知若只說方言，不用啟示，或知識，或預言，或教訓來講解就對衆人沒有益處（林前14:6）。先知講道是為信的人做證據。要切慕做先知講道。
說方言（different kinds of tongues7） 按照聖經教導，說方言之能終必停止（林前13:8）（本經文所題到的要多深入解釋）。英文的方言是複數，所以方言不衹是一種8。按照方語的種類分有: 1，萬人的方言（林前13:1） 2，天使的言語（林前13:1） 保羅盼望信徒人人都可以說方言（林前14:5），由此可以推知這個恩賜是聖靈賜給每一位信徒的恩賜。說方言之能終必停止。
3. Life Application Study Bible（NIV）, 2081.
4. Life Application Study Bible（NIV）, 2050.
5. Life Application Study Bible（NIV）, 2050.
6. Life Application Study Bible（NIV）, 2081.
7. Life Application Study Bible（NIV）, 2081.
8. The NIV Study Bible, (Grand Rapids, MI : Zondervan Publishing House, 1995), 1752.
4
如果按照說方言的内容分類有:
1，信息方言（會衆聚集的時候）
2，禱告方言（自己獨處的時候）
下面按照第二種分類方法來討論說方言的恩賜。
信息方言
1，會衆聚集的時候，說的方言也可能是地上萬人的方言，也有可能是天使的言語。 2，說方言不是為信的人做證據，乃是為不信的人（林前14:22a）
說方言的若不翻出來，使教會被造就，那做先知講道的就比他強（林前14:5） 那說方言的，就當求著能翻出來（林前14:13）。若沒有人翻，就當在會中閉口（林前14:28）。全教會聚在一處的時候，不要都說方言（林前14:23），若有說方言的，只好兩個人，至多三個人，且要輪流著說，也要一個人翻出來（林前14:27）。要切慕做先知講道，也不要禁止說方言（林前14:39）
禱告方言 禱告的方言不是對人說，乃是對神說，因為沒有人聽出來，然而人在心靈裡卻是講說各樣的奧祕（林前14:2）。這種方言指用天使的言語向神禱告。保羅說這種方言比衆人都多（林前14:18），但他在教會中寧願說造就教會的話（林前14:19）。 說這種方言人的是造就自己的靈（林前14:4a），這種禱告的言語是靈在禱告，悟性沒有果效（林前14:14）。信徒既要用靈禱告也要用悟性禱告（林前14:15a）。
翻方言 「翻方言」是「聖靈的恩賜」（林前12:10），適用於會眾聚會時使用而且是造就教會的恩賜（林前14:5）
「說方言講道」時倘若沒有「翻方言」的人，就當閉口不言。「說方言講道」和「翻方言」是同時出現的恩賜（林前14:13）
教導
當專一
勸化
當專一
C. 異能的恩賜
異能的恩賜包括醫病（這裏所説的醫病不是用人類的醫療知識和設備，乃是神用超自然的方式醫治人類用現有的科技手段和知識無法醫治的疾病，或者使患者立刻恢復健康的醫治方式），行異能，和辨別諸靈的恩賜。這次的必讀經文中沒有提及行異能的恩賜，所以這裏不做介紹。
恩賜
簡單説明
運用及原則
醫病
醫病的恩賜即可以醫治身體又可以醫治心靈。
5
gifts of healing9
英文NIV用的是gifts of healing，恩赐用的是复数。
行異能
辨別諸靈
辨别諸靈的恩赐中的靈是复数。
先知的靈原是順服先知的（林前14:32）
D. 工作的恩賜
工作的恩賜包括信心，執事（serve），施捨，治理和憐憫的恩賜。下面列表逐一簡單説明。
恩賜
彰顯的情形
運用及原則
信心 可以用於先知的講論 照著信心的程度說預言（羅12:6） 用於正確的審查評價自己 信徒當照著神分給個人信心的大小看的合乎中道（羅12:3）
執事Serve10
執事的NIV英文是serve，意思是服事。 應當專一
施捨
當誠實
治理Lead11
治理的NIV英文是領導的意思 當殷勤
憐憫
當甘心
\subsubsection{恩賜的擁有人（林前12:12-13，弗4:1-10，羅12:3-8）}
恩賜雖然不同，但只要是信徒，就是恩賜的擁有者。保羅用身體和肢體比喻信徒和基督的關係（羅12:4-6），每一個信徒都是一個肢體，雖然以前的身份不同，如不同的種族或社會的地位，身份，但只要受洗歸入基督，就是從一位聖靈的感動（林前12:13），聖靈就把恩賜放在不同的肢體中。而且個人蒙恩都是照基督所量給各人的恩賜（弗4:7）說的也是一樣的道理，即每個蒙恩的人，都有神給的恩賜在身上。所以，每一位信主的信徒，都是恩賜的擁有者，不能再說我們沒有恩賜了。
V. 各恩賜的重要（林前12:14-18，弗4:11-14，羅12:3-8）
A. 每個恩賜都不可缺
人身上有許多的肢體，如果缺了任何一個，叫殘疾。而且每個肢體之間都是互為需要，缺了眼睛，即使有腿走路也容易摔跤，所以保羅把每個信徒都比作一個肢體，身體需要每一個肢體才完全，才是健康的，而每個肢體身上都有聖靈賜給的恩賜，因此，可以說每個恩賜對教會來說也都是不可或缺的。
B. 每個恩賜盡其職能
雖然恩賜各有不同，如使徒，先知，教師，傳福音的等等，但每個恩賜都不可缺少，不僅不可缺，而且每個恩賜都有不同的職能，如傳福音的就是要按大使命去傳福音，使徒早期的工作是到處建立教會，或像保羅把聖靈給他的啟示寫出來給各地的教會，教師把聖經的真理教導清清楚楚，使信徒信仰的根基穩固等等，每個恩賜的職能都要被發揮出來，即『各盡其職』（弗4:12），才能建立基督的身體，即使有異教之風，也不致被搖動。
C. 每個恩賜互相聯絡
雖然恩賜各有不同，但不是每個恩賜互不相關，也不是有哪一個恩賜大到不需要其他的
9. Life Application Study Bible（NIV）, 2081.
10. Life Application Study Bible（NIV）, 2082.
11. Life Application Study Bible（NIV）, 2082.
6
恩賜，就如身上的肢體一樣，各自有不同的用途，恩賜也是一樣，雖然每個信徒都要在他的恩賜上要專一（作好），但要互相聯絡在一起，互補才能建造健康的身體，不能輕看他人的恩賜或高看自己的恩賜，因為所有的事都是神賜給的，沒有一樣是我們自己靠努力或能力得來的，所以要看得合乎中道才好（羅12:3b）。
D. 恩賜不可能脫離教會以外存在（林前 12:15）
就如同人的耳朵長在身體上，就是健康的，可以正常工作；但如果離開身體，很快就失去生命，如果我們脫離了教會，我們的恩賜就會退化，甚至失去功能。
VI. 面對自己的恩賜（林前12:19-24，弗 4:15-16，羅 2:3-8）
A. 我不一定與別人有一樣的恩賜
保羅告訴我們每個人都有恩賜，就好比身體只有一個，但肢體是多的（林前12:19-21）而且用途不一樣（羅12:4），所以我的恩賜與別人的恩賜可能不一樣，教會的信徒也不可能都有一樣的恩賜。
B. 我有的恩賜與別人的一樣重要
每個肢體都有各自的功用，不能說一個肢體不需要另個肢體（林前12:19-21），所以神給我的恩賜，也給別的信徒恩賜，我不能因自己的恩賜就覺得自己高於他人，其他信徒也是如此，都要看到合乎中道（羅12:3），因為恩賜都是神賜給，不是自己努力得來的。
C. 我認為不重要的恩賜神不輕看
保羅說人以為軟弱的肢體，更不可少，人看為不體面的，神卻加上體面，讓它更俊美（林前12:23-24）同樣，恩賜沒有哪一個是重要，哪一個是不重要的，是神給的，對教會來說都一樣重要，我們不要輕看自己的恩賜，也不要輕看其他信徒的恩賜。
D. 我的恩賜與別人的要互相聯絡
每個是肢體都要連於元首基督，而且聯絡的合適，各自發揮自己的功用，才能叫身體漸漸增長（弗4:15-16a），所以我們各自的恩賜，不是單獨的發揮作用，而是與其他肢體的恩賜一起，各按各職，而且要彼此互助的，彼此之間的媒介就是愛，這樣才能真正發揮作用，不但建造教會，也會在愛中建立自己（弗4:16b）。
VII. 恩賜的配搭（林前12:19-27，弗4:1-6，羅12:3-8）
一個肢體不能成為身體，身體也不可能只有一個或都是一樣的肢體，因為肢體的不同，而且互相的配搭才能叫身體活動自如，才能健康。同樣聖靈所賜的恩賜也是不同的，多樣的，因此也需要彼此的配搭才可以建立基督的身體。恩賜的配搭原則包括:
A. 彼此需要（林前12:21）
身體上有許多的肢體，每個肢體都是運用合適，身體才能正常的活動，手沒有眼睛，就不能準確的拿到需要的東西，眼睛看見了，沒有手也是拿不到東西，可見肢體互相的需要才可以完成一個動作，一項任務。這種互相的需要在恩賜上也是一樣。在神的教會裡，一個恩賜不能使教會被建造，比如有傳福音的恩賜的人，把福音傳給人，接下來就需要有人去關心新的信徒，也需要有人去教導他們聖經的真理才能在真理上成長，如果軟弱了還需要有人去勸勉，有的信徒生活有需要了，還要有人願意奉獻所需等，這就是恩賜上的彼此需要。
B. 彼此看重（林前12:22-24，弗4:1-2a）
身體上的肢體，有的很容易受傷，比如眼睛，看起來好像很軟弱，但如果沒有眼睛行動就困難。所以，身體上的肢體缺一不可。同樣，恩賜也是如此，聖靈賜給教會信徒的恩賜，每一個都有不同的功用，有人恩賜多一些，有人少一些，但這是神分配的，所以信徒之間不可以在恩賜上彼此輕看，過度看重自己或每個人的恩賜，忽略其他人的恩賜；反而要彼此看重，個人要謙卑、溫柔、忍耐，互相寬容，這樣才與蒙召的恩相稱。
7
C. 彼此相顧（林前12:25）
身體上的肢體互相聯絡，才可以有身體的和諧的動作。當一隻手拿東西要掉落時，眼睛看見，大腦反應，另一隻手就會立刻幫助接住，腿也要支撐，腰也要用力，而不是眼睛看見卻幸災樂禍不管，另一隻手也躲開。恩賜的服事也不能分開，有的恩賜服事會有明顯的果效，如有傳福音恩賜的信徒，他一講，有人就信主，其他人的反應是與他一起歡喜，還是心生嫉妒？還是一起去關心新的信徒，幫助他成長呢？或者有人恩賜不明顯，服事不得力，很灰心，就如受苦的肢體一樣，其他信徒就應當去陪伴，看如何去幫助他。這就是恩賜服事上的彼此相顧，而不是彼此嫉妒。
D. 彼此合一（弗4:2）
不論信徒有什麼樣的恩賜，首先要讓自己的恩賜發揮出來，並且要專一（羅12:6-8），照著神分給各自的恩賜在教會裡作好，比如教導的就專一教導，就是把教導的工作作好。憐憫人的，要甘心的作等。雖然每個信徒要專一作自己的服事，但心裡當清楚，神是一位，主是一位，聖靈也是一位，所以各自發揮各自的恩賜，都是要服事同一位主，恩賜來與同一位聖靈，因此，恩賜的服事合一更重要，無論什麼樣的恩賜，什麼樣的服事，什麼樣的職分，都要保守聖靈所賜合而為一的心。
VIII. 恩賜的配合（林前12:27-31，弗4:1-10，羅12:3-8）
恩賜是聖靈所賜，分配在不同的肢體身上，恩賜有多種（弗4:8b），分給各人的也不同，在教會裡不都是使徒，也不會都是教師（林前12:29-30），所以，恩賜的服事需要彼此配搭，互相配合才能發揮所有肢體不同的用處（羅12:4）。而且在神的教會裡，神的心意是各司其職，即每個肢體都要參與服事，把所有恩賜都充分發揮使用出來，才能建立基督的身體。比如说方言和翻方言的恩賜是彼此配合最明顯的實例，有人有說方言的恩賜，如果在聚會時，他說方言或用方言禱告，沒有會翻譯方言的恩賜的信徒在場，別人聽不懂，就不知道他在說什麼，禱告感謝後，別人也不會說阿門（羅14:16）。反之，如果沒有人會說方言，翻方言的恩賜也不可能體現出來，因此，教會裡恩賜配合好才能真正體現其意義和用途。
IX. 在愛中使用恩賜（最妙之道）（林前13章）
A. 愛比恩賜更重要的原因（13:1-3）
保羅講了各樣的恩賜，但到最後，他告訴教會信徒要渴慕更大的恩賜，就是愛。如果沒有愛，恩賜會成為什麼呢？
1. 沒有愛的恩賜是鳴的鑼響的鈸（林前13:1）
當我們信主後，主的愛就住在我們裡面。聖靈的恩賜主要體現在外在的表現，比如言語方面，能力方面，或具體的工作方面，如果只有注重外在的表現，只為了完成一項工作或一項服事，忽略了內在的因素，比如謙卑，溫柔，忍耐，和平，寬容等，不顧被一起服事或被服事者的內心感受，不會造就人，反而會傷害人。就如鳴的鑼響的鈸，外面聲音大而已。
2. 沒有愛的恩賜是與我們無益的（林前13:2-3）
如果一個人有信心可以挪山，或者把所有一切周濟窮人，不是出於愛的緣故，不是為了愛神，不是為了愛服事的對象，只是顯出了自己的恩賜，就沒有任何的益處。
B. 愛的重要性（13:4-8）
1. 對待別人的（林前13:4）
我們對愛的定義，是看別人對我們有沒有愛的表現，但保羅這裡重新給我們定義了愛-那就是我們對別人有沒有付出愛。首先說愛是恆久忍耐和恩慈，在恩賜的服事中，如果對同工有忍耐，就不會只是看重效率。如果給人傳福音有忍耐，就會注重對福音對象的關心，而不是他能不能立刻信主。如果傳福音的對象反對，甚至用言语攻
8
擊我們時，更需要忍耐，甚至還要對他跟有恩慈，如果我们也用言语回击对方，我们就没有好的见证。愛是不嫉妒，嫉妒是看到他人比自己強，或看到他人有的比自己多，或看到他人比自己好而產生的內心的活動，在恩賜服事上不可有嫉妒之心，反而要與喜樂的人同樂，當其他信徒在恩賜服事上有果效，甚至比自己作的好，要與他一同快樂。
2. 不要作的事（林前13:5-6a）
愛是不嫉妒，不自誇，如果自己在恩賜服事上有果效，卻要有『是不自誇，不張狂』才好。恩賜的服事也不要期望為自己地回報，不要求自己的益處，為他人著想，別人作錯了，也不要記在心裡，不計算人的惡。
3. 要作的事（林前13:6-7）
凡事包容，包容是因為別人作的不夠好，我們可以接納，或者別人作了傷害我們的事，我們也可以忍耐，這對一起服事的同工尤其重要。凡事相信，不是相信人，而是相信神的心意是好的；凡事盼望也是在恩賜服事上短時看不到果效時，還對神有信心，相信神有一天會成就；凡事忍耐，要作到凡事包容，凡事相信，凡事盼望都需要有忍耐的心。
4. 愛永存（林前13:8）:當主再來時，恩賜會終止，結束，但愛卻會存到永遠。
C. 恩賜與愛的影響力（林前13:8-13）
恩賜是為要建造教會，成全聖徒而用，當主耶穌再來時，我們與主面對面，現在不明白的，模糊的，無論是知識上的，或有關天國的事，就都清楚了，全都知道了，所以，先知講道的恩賜不再需要，因為我們不需要有人來教導了，到那時所知道的，就如同主知道我們一樣。因此，恩賜的服事是暫時的，如果只注重恩賜，沒有愛，我們所作的一切就不會存到永遠。
X. 我們對待恩賜的態度
A. 愛是恩賜服事的根基（所有的恩賜在使用時必須有愛，應當在愛中切慕恩賜（林前14:1）
恩賜服事中，要把愛放在首位，愛是恩賜服事的出發點，愛要貫穿於恩賜服事的整個過程。
B. 不要高看自己
恩賜是聖靈按己意賜給信徒的（林前12:8，弗4:7），不是信徒自己努力得來的，所以，無論有多少恩賜，用恩賜服事有多大的果效，都不要驕傲，是聖靈借我們來作的工作，要把榮耀歸給神。
C. 不要輕看別人
恩賜是用來建造教會，成全聖徒（弗4:12），這就要求恩賜的服事是信徒之間彼此的互動，而不是一個人單獨的工作，沒有其他同工（恩賜）的配合，不能完成教會群體的建造，就如保羅說的，肢體之間不能說我不需要其他的肢體。
D. 要切慕屬靈的恩賜，並要求多得造就教會的恩賜（林前14:1,12）
恩賜有多種，保羅教導信徒在追去愛的同時，也要切慕屬靈的恩賜，更有羨慕的是作先知講道的恩賜，因為先知講道的恩賜是造就教會的，而有的恩賜是造就自己的，如說方言的（林前14:1-4）。
E. 願意與他人搭配（林前12:27，29-30） 恩賜的服事好比身體上的肢體，需要彼此的配搭，才能起到建造教會的作用。所以在信徒清楚自己的恩賜同時，也要在服事中主動尋找有其他恩賜的同工一起配搭，教會是個群體，不是個體，群體的服事中要學會與他人的配合。
F. 各別恩賜需要特別的態度 恩賜服事需要與其他同工一起配搭，但不代表可以忽略自己恩賜服事的重要，個體服事
9
的重要性表現在不同的恩賜，有的需要特別的態度來使用這些恩賜。
1. 信心的態度（羅12:6）
說預言的的恩賜，在使用時，要照著信心的程度說預言，不能誇大或減小。
2. 專一的態度（羅12:7-8a） 需要有專一態度的恩賜包括執事，教導和勸化，如果一位同工有講道的恩賜，就讓他專一教導，或者作執事，就專一作執事，作勸化的，就專一作勸化，在這些恩賜服事時盡量不要被其他事物打擾。
3. 誠實的態度（羅12:8）
需要誠實的服事態度的恩賜是施捨的恩賜，在施捨時要誠實，或叫大方。
4. 殷勤的態度（羅12:8）
有治理恩賜的，在治理時要殷勤，不可懶惰。
5. 甘心的態度（羅12:8）
有憐憫的恩賜的人，在服事時要甘心樂意，不要表現出被迫在作。
G. 不要排斥恩賜（林前14:1）
保羅教導說聖靈顯在各人身上是叫人的益處，信徒要切慕屬靈的恩賜，服事中會有失敗，挫折，軟弱，但不應以此就拒絕恩賜，排斥恩賜，因為恩賜的服事會讓基督的身體漸漸增長，並且可以在愛中建立自己（弗4:16）。
H. 不輕忽恩賜（林前12:23） 保羅用身體上人以為軟如的肢體更是不可少的，他在告訴我們恩賜中我們以為小的，果效不明顯的，或者在我們眼中看不上的，反而是更重要的，所以，我們不要輕忽任何一個恩賜。
I. 不強求別人與我一樣（林前12:4-6）
聖靈所賜的恩賜是不同的，當我們作一項事工時需要同工，有時我們就期望每個人都有跟我們一樣的恩賜就好了，甚至強迫別人作不是他恩賜所擅長的服事。
J. 同苦同樂的心，別人得榮耀了我們不要嫉妒（林前12:26）
保羅講到用愛來服事時，說到愛是不嫉妒，這在恩賜服事上也很重要，當其他同工服事有效，有成就時，我們的心態是什麼？是嫉妒，為什麼神給他那麼多恩賜不給我？還是與他一同高興快樂，看到神家的事工有果效，我們一起快樂，一起得榮耀？所以在恩賜服事中要有同苦同樂的態度才對。
K. 不要禁止恩赐（林前14:39）
所以恩賜都是聖靈所賜，都是為成全聖徒，建造教會，所以不要有禁止某個恩賜的心。保羅說要切慕先知講道，但也不要禁止說方言。
L. 順服的態度（林前14:29-32）
在恩賜服事中也要有彼此順服的態度，如果先知講道的，旁邊的人得了啟示，先說話的就要閉口，不但要順服聖靈的引導，也要彼此順服。
XI. 我的學習
A. 透過學習，我糾正了一些以前對恩賜認識不正確的地方。比如我一直以爲所有的恩賜都不是每個人可以擁有的。但根據聖經的教導，人人都可以擁有先知講道，說方言和智慧言語的恩賜。
B. 以前對聖靈的恩賜認識不全面的地方，透過這次學習有了更深入的瞭解，比如我以前把先知的恩賜認爲是“可以說出未來的事情”，其實這衹是先知恩賜的一部分功能。其實先知的恩賜還包括：在教會的講道。
C. 我發現恩賜可以通過操練而成長。
1. 純粹從神那裏所得來的，比如行異能和醫病的恩賜。這些恩賜不是天生的像有的人
10
天生記憶力好那樣，也不是後天透過學習或者刻苦操練可以得到。
2. 有的恩賜是從信主后通過操練而成長的。比如信心和先知的恩賜。保羅說做先知的要照著信心的程度說預言，可見先知的恩賜是與信心的操練同步成長的。
D. 我受洗后一直都在保守的福音教會成長，從來沒有如此仔細的查考過聖經中對於聖靈的恩賜的描述和記載，因此對聖靈的恩賜有些排斥，怕走偏了。這次通過學習，聖靈的恩賜是天父賜給教會的禮物，是用來幫助教會成長的恩典。我盼望可以可到更多的建造教會的屬靈的恩賜。
參考書目
1. The NIV Study Bible. Grand Rapids, MI: ZondervanPublishingHouse, 1995
2. Life Application Study Bible（NIV）. Wheaton, IL: Tyndale House Publishes, Inc., 1991.

\subsection{基本真理}


\subsubsection{神的奧秘}
\begin{itemize}
\itemsep-.25em 
\item 奧秘的意思 
\begin{itemize}
\itemsep-.25em 

\item 神在萬世以前\underline{預定}並\underline{隱藏}的事情。
\begin{itemize}
\itemsep-.25em 
\item 「我們講的，乃是從前所隱藏、神奧秘的智慧，就是神在萬世以前預定使我們得榮耀的。 這智慧，世上有權有位的人沒有一個知道的，他們若知道，就不把榮耀的主釘在十字架上了。  如經上所記：「神為愛他的人所預備的，是眼睛未曾看見，耳朵未曾聽見，人心也未曾想到的。」 林前2:7-9
\end{itemize} 

\item 在\underline{新約書信前}是沒有人知道的事情。
\begin{itemize}
\itemsep-.25em
\item 「這奧秘在以前的世代沒有叫人知道」——弗3:5a
\end{itemize}
\item 耶穌來到世上向人類啟示自己
\begin{itemize}
\itemsep-.25em
\item 這是要應驗先知的話說：「我要開口用比喻，把創世以來所隱藏的事發明出來。」( 太13:35,詩78:2)
\end{itemize} 
\item 現在藉著\underline{聖靈啟示}給\underline{使徒}和\underline{先知}的事情。
\begin{itemize}
\itemsep-.25em 
\item 「像如今藉著聖靈啟示他的聖使徒和先知一樣。」——弗3:5b 
\item 都是照他自己所預定的美意，叫我們知道他旨意的奧祕——弗1:9
\end{itemize}
 
\item 神願意將這奧秘啟示給一切信祂的人
\begin{itemize}
\itemsep-.25em
\item 耶穌回答說：「因為天國的奧祕只叫你們知道，不叫他們知道」。(太13:11)

\end{itemize}

\end{itemize}

\item 隱藏的原因 
\begin{itemize}
\itemsep-.34em 
\item 人的智慧有限。 「除了在人裡頭的靈，誰知道人的事；像這樣，除了神的靈，也沒有人知道神的事」——林前2:11 
\item 先讓律法來臨,令人知罪，在律法仍有功用的時刻，關乎救恩的奧秘就隱藏了。 「但這因信得救的理還未來以先，我們被看守在律法之下，直圈到那將來的真道顯明出來。這樣，律法是我們訓蒙的師傅，引我們到基督那裡，使我們因信稱義。但這因信得救的理既然來到，我們從此就不在師傅的手下了」——加3:23-25 
\end{itemize}
\item 神的奧秘就是基督
\begin{itemize}
\itemsep-.25em 
\item 我願意你們曉得，我為你們和老底嘉人，並一切沒有與我親自見面的人，是何等地盡心竭力，要叫他們的心得安慰，因愛心互相聯絡，以致豐豐足足在悟性中有充足的信心，使他們真知\underline{神的奧祕，就是基督}，所積蓄的一切\underline{智慧知識都在他裡面藏著}。(西2:1-3)
\item 這道理就是歷世歷代所隱藏的奧祕，但如今向他的聖徒顯明了。  神願意叫他們知道，這奧祕\underline{在外邦人中}有何等豐盛的榮耀，就是\underline{基督在你們心裡成了有榮耀的盼望}。(西1:26,27)
\item 你們從前不順服神，如今\underline{因他們的不順服，你們倒蒙了憐恤}。  (羅16:30)
\item 這樣，他們也是不順服，叫他們因著施給你們的憐恤，現在也就蒙憐恤。弟兄們，我不願意你們不知道這奧祕，恐怕你們自以為聰明，就是：以色列人有幾分是硬心的，等到\underline{外邦人的數目添滿}了， 於是\underline{以色列全家都要得救}。如經上所記：「必有一位救主從錫安出來，要消除雅各家的一切罪惡。」又說：「我除去他們罪的時候，這就是我與他們所立的約。」 (羅16:31,25-28)
\item 因為神將眾人都圈在不順服之中，\underline{特意要憐恤眾人}。(羅
16:32)
\end{itemize}




\item 神的心意是把這個奧秘藉著教會顯明出來
\begin{itemize}
\itemsep-.25em 
\item 要照所安排的，在日期滿足的時候，使天上、地上一切所有的，都在基督裡面同歸於一。我們也在他裡面得了基業，這原是那位隨己意行做萬事的，照著他旨意所預定的，叫他的榮耀從我們這首先在基督裡有盼望的人，可以得著稱讚。 你們既聽見真理的道，就是那叫你們得救的福音，也信了基督，既然信他，就受了所應許的聖靈為印記。這聖靈是我們得基業的憑據 ,直等到神之民 被贖，使他的榮耀得著稱讚。(弗1:10-14)
\item 用啟示使我知道福音的奧祕，正如我以前略略寫過的。你們念了，就能曉得我深知基督的奧祕,這奧祕在以前的世代沒有叫人知道，像如今藉著聖靈啟示他的聖使徒和先知一樣。這奧祕就是外邦人在基督耶穌裡，藉著福音，得以同為後嗣，同為一體，同蒙應許。又使眾人都明白，這歷代以來隱藏在創造萬物之神裡的奧祕是如何安排的，為要藉著教會，使天上執政的、掌權的現在得知神百般的智慧。這是照神從萬世以前在我們主基督耶穌裡所定的旨意。(弗3:3-6,9-11)
\item 也為我祈求，使我得著口才，能以放膽開口講明福音的奧祕——我為這福音的奧祕做了戴鎖鏈的使者——並使我照著當盡的本分放膽講論。(弗6:19-20)
\end{itemize}
\item 保羅做了傳講這奧秘事的執事
\begin{itemize}
\itemsep-.34em
\item 你們要恆切禱告，在此警醒、感恩。也要為我們禱告，求神給我們開傳道的門，能以講基督的奧祕——我為此被捆鎖——叫我按著所該說的話將這奧祕發明出來。(西4:2-4) 
\item 唯有神能照我所傳的福音和所講的耶穌基督，並照永古隱藏不言的奧祕，堅固你們的心。這奧祕如今顯明出來，而且按著永生神的命，藉眾先知的書指示萬國的民，使他們信服真道。(羅16:25-27)


\item 只有神藉著聖靈向我們顯明了。因為聖靈參透萬事，就是神深奧的事也參透了。  除了在人裡頭的靈，誰知道人的事？像這樣，除了神的靈，也沒有人知道神的事。  我們所領受的，並不是世上的靈，乃是從神來的靈，叫我們能知道神開恩賜給我們的事。  並且我們講說這些事，不是用人智慧所指教的言語，乃是用聖靈所指教的言語，將屬靈的話解釋屬靈的事 。 然而，屬血氣的人不領會神聖靈的事，反倒以為愚拙；並且不能知道，因為這些事唯有屬靈的人才能看透。  屬靈的人能看透萬事，卻沒有一人能看透了他。  「誰曾知道主的心，去教導他呢？」但我們是有基督的心了。(林前2:10-15)
\item 人應當以我們為基督的執事，為神奧祕事的管家。(林前4:1)


\end{itemize}
\end{itemize}

\subsubsection{因信稱義}

\subsubsection{律法}

\subsubsection{割礼}

\subsubsection{信心}

\subsubsection{復活}

\subsection{教會生活}
\subsubsection{信徒生活}

\subsubsection{教會問題處理}

\subsubsection{教會治理智慧}

\subsubsection{如何说感谢的话}
感谢神，教会，人
腓利比书1:3

\subsection{救恩}
林多前書 3
我照神所給我的恩，好像一個聰明的工頭立好了根基，有別人在上面建造，只是各人要謹慎怎樣在上面建造。 11 因為那已經立好的根基就是耶穌基督，此外沒有人能立別的根基。 12 若有人用金、銀、寶石、草、木、禾秸在這根基上建造， 13 各人的工程必然顯露；因為那日子要將它表明出來，有火發現，這火要試驗各人的工程怎樣。 14 人在那根基上所建造的工程若存得住，他就要得賞賜； 15 人的工程若被燒了，他就要受虧損，自己卻要得救，雖然得救，乃像從火裡經過的一樣。


救赎是从危险或苦难中解脱出来。得救是解脱与保护。这个词包含的意义有胜利，健康，保守。圣经有时候用saved或salvation这两个词指暂时的、身体的解脱，比如保罗从狱中得救（腓立比书1:19）。

‘salvation’这个词更多的指永恒的，精神上的解脱。当保罗告诉腓立比的囚犯该怎样做才能得救时，他指 的是这个囚犯永久的命运（使徒行传16:30-31）。耶稣把得救等同于进入神的国（马太福音19:24-25）。

我们从何得救？基督教关于救赎的教义中，我们从“怒”得救，也就是说，从神对罪的审判得救（罗马书5:9；帖撒罗尼迦前书5:9）。我们的罪将我们与神分隔，罪的结局就是死（罗马书6:23）。圣经的救恩是指我们从罪的恶果中解脱出来，所以包括了罪的移除。

谁来拯救？只有神可以消除罪恶，并将我们从罪的刑罚中拯救出来（提摩太后书1:9；提多书3:5）。

神怎样救我们？在基督教救赎的教义中，上帝通过基督救了我们（约翰福音3:17）。具体地说，正是耶稣死在十字架上以及随后复活，实现了对我们的救恩（罗马书5:10；以弗所书1:7）。圣经明示，救恩是神恩惠的、不应得的恩赐（以弗所书2:5，8），并只能藉着信耶稣基督得到（使徒行传4:12）。

我们如何得救？藉由信。首先，我们必须听到福音 — 听到耶稣死亡和复活的好消息（以弗所书1:13）。然后，我们必须相信 — 完全信主耶稣（罗马书1:16）。包括悔改，对罪和基督认识上的转变（使徒行传3:19），求告主名（罗马书10:9-10，13）。

基督教救恩教义的定义是“藉着神的恩，赐给那些信神并接受忏悔的条件，且信主耶稣的人从罪永久的惩罚中解脱出来。”得救只有在耶稣里（约翰福音14:6，使徒行传4:12），而且只能靠神提供、担保和守护



\subsection{國度觀}

